\documentclass{sprawozdanie-agh}

\let\lll\undefined

\usepackage{amssymb}
\usepackage[utf8]{inputenc}
\usepackage{matlab-prettifier}
\usepackage{siunitx}

\makeatletter

% ========= Początek Dokumentu ========= 

\begin{document}


% ========= Strona tytułowa ========= 
\przedmiot{Teoria Sterowania II: Ćwiczenia Laboratoryjne}
\tytul{Wytyczne dla Ćwiczenia 1 }
\podtytul{Portrety Fazowe Systemów~Liniowych}
\kierunek{Automatyke i~Robotyka}
\autor{Dariusz Cieślar}
\data{Kraków, \today}

\stronatytulowa{}

% ========= Treść Właściwa =========

\section{Przebieg ćwiczenia}

\begin{description}

\item[Krok 1] Przygotuj 9 macierzy $A$ (rozmiar 2x2) spełniające wszystkie przypadki~określone w~tabeli w~skrypcie. Jeśli~istnieje możliwość wyboru wartości własnych dodatnich lub ujemnych, proszę wybrać ujemne.

\item[Krok 2] Zaimplementuj funkcję w~MATLABie pozwalającą na symulację i~rejestrację przebiegów~czasowych dla obu zmiennych stanu. Funkcja ta powinna umożliwiać zmianę macierzy $A$ oraz warunków~początkowych, najlepiej w~postaci~argumentów. 

\item[Krok 3] Przeprowadź symulację dla wybranego przypadku (wybranej macierzy $A$ oraz wybranego warunku początkowego), zweryfikuj poprawność symulacji i~przygotuj skrypt lub funkcję pozwalającą na:
\begin{itemize}
    \item wyrysowanie trajektorii~stanu na płaszczyźnie stanów,
    \item wyrysowanie podprzestrzeni~niezmienniczych na płaszczyźnie stanów,
    \item wyrysowanie wektorów~własnych/głównych,
    \item oznaczenie punktów~równowagi,
    \item opcjonalnie: zaznacz pochodne stanu dla siatki~punktów~w~przestrzeni~stanów~(np. za pomocą funkcji~\textit{quiver}),
    \item wyrysowanie przebiegów~czasowych zmiennych stanu na osobym wykresie.
\end{itemize}

\item[Krok 4] Przeprowadź symulacje dla szeregu warunków~początkowych (np. rozłożonych równomiernie po obwodzie koła). Wykorzystaj funkcje wykreślania wyników~z kroku 3, by wykreślić portety fazowe, tzn. wykreślic zbiorczy wykres trajektorii stanów~dla wszystkich badanych warunków~początkowych na płaszczyźnie stanów. Do zbioru warunków początkowych dodaj punkty leżące w~podprzestrzeniach niezmienniczych.

\item[Krok 5] Wykonaj przekształcenie macierzy $A$ z kroku 3 do postaci~kanonicznej i~powtórz krok 4. Jeżeli~macierz $A$ była już w~postaci~kanonicznej - przekształć ją do innej postaci~i~powtórz krok 4.

\item[Krok 6] Dla macierzy wykorzystywanej w~kroku 3, zamień znaki~wartości~własnych (ujemne na dodatnie) i~powtórzy krok 4.

\item[Krok 7] Przeprowadż symulacje opisane w~kroku 4 dla wszystkich 9 macierzy przygotowanych w~kroku 1. Zweryfikuj, że uzyskane portrety fazowe wyczerpują przypadki opisane w~skrypcie. Jeśli zaobserwujesz odstępstwa, ponownie przyjrzyj się swojemu wyborowi przykładów~macierzy $A$ i~skoryguj błędy.

\end{description}

\textbf{Zadanie dodatkowe}
Zaproponuj fizyczną realizację układu drugiego rzędu tak, aby modyfikowanie wartości elementów~tego układu pozwoliło na uzyskanie co najmniej 5 typów portretów fazowych.
Wykreśl schemat oraz sformułuj model fizyczny tego układu. Wyznacz przykładowe wartości parametrów tego modelu dla każdego z~typów~portretów fazowych.

\section{Opracowanie wyników}

Przygotuj sprawozdanie z~ćwiczenia opisujące cel, przebieg i~obserwacje. W~części opisującej przebieg ćwiczenia proszę przedstawić w~sumie jedenaście przypadków:
\begin{enumerate}
    \item macierz z~pierwotnego przykładu (kroki 2, 3 i~4),
    \item przekształcenie macierzy z pierwotnego przykładu do postaci kanonicznej,
    \item zmiana znaków wartości~własnych z~ujemnych na dodatnie dla pierwotnego przykładu,
    \item osiem macierzy odmiennych od pierwotnego przykładu w~celu wyczerpania wszystkich dziewięciu typów~portretów~opisanych w~skrypcie.
\end{enumerate}

Dla każdego z przypadków~należy przedstawić:
\begin{itemize}
    \item macierz stanu, jej wartości~własne i~wektory własne/główne. 
    \item porter fazowy zawierający:
    \begin {enumerate}
        \item trajektorie dla co najmniej 10 warunków~początkowych, w~tym warunków leżących w~podprzestrzeniach niezmienniczych,
        \item wektory własne/główne,
        \item punkty równowagi,
        \item przestrzenie niezmiennicze,
        \item wyróżniony wykres trajektorii dla jednego z warunków~początkowych,
    \end {enumerate}
    \item wykres z~przebiegami czasowymi~zmiennych stany dla warunku początkowego wyróżnionego na portrecie fazowym.
\end{itemize}

\textbf{Zadanie dodatkowe}
Proszę przedstawić rozwiązanie zadania dodatkowego w~postaci:
\begin{itemize}
    \item schematu proponowanego układu,
    \item równań różniczkowych opisujących model fizyczny tego układu,
    \item równania stanu z~postacią macierzy stanu w~formie funkcji~parametrów~układu,
    \item tabelę pokazującą parametry elementów~układu, postać macierzy $A$, wartości~własne dla tych parametrów~oraz portety fazowe.
\end{itemize}

\end{document}
